% !TeX root = ../papers/nc-p3.tex
\section{非対称な四窓の滑らかな三文字切断と二つの鉛筆}
\label{sec:finite-determinantal-marked-slice}
\begin{theorem}\label{thm:finite-determinantal-marked-slice}
有限な四辺を持つ標準有限Frobenius窓に $x\in\mathbb P(V_0^*)$ を印し、
$t=w(x)\in V_1\otimes V_2\otimes V_3$ の三つの四次行列式曲面が滑らかとする。
中央の曲面 $T\subset\mathbb P(V_2^*)$ はK3である。順序付き三偏極を
$l,h,r$ と書けば
\[
 l^2=h^2=r^2=4,\quad lh=hr=6,\quad l+r=3h .
\]
$A=l+h, B=h+r$ とすると、元の長さ20の点グラフは $T$ 上で
$|A|$ と $|B|$ のそれぞれの鉛筆の基底スキームとなる。
一方の鉛筆を指定すれば他方は一意である。

逆に三文字 $t$ を固定すると、これを指定切断とする全ての窓を次の有限方程式で復元できる。
$\Lambda\in\mathrm{Gr}(2,H^0(T,A))=\mathrm{Gr}(2,12)$ の基底スキーム $X$ が零次元で、
$H^0(I_{X/T}(B))$ が二次元かつその基底スキームも零次元となることを課す。
二鉛筆をそれぞれ双線型式へ引き戻して得た六空間を $R_1,R_2$ とする。
三次重なりが四次元であることと通常の有限Frobenius・有限辺検査を課し、
その重なりに $t$ を印す。これでちょうど指定した候補を得る。
印付き添字保持同型は順序を保つ $t$ の三基底変更と鉛筆の対応と同値である。
\end{theorem}
\begin{proof}
$t$ を中央点の四行四列線型行列と読む。滑らかな行列式上では階数は全点で3であり、左右核点が定まる。
他の二位置の行列式でも同じため、この二つの核写像は四次曲面どうしの同型である。
隣り合う二射影のグラフは四つの $(1,1)$ 式の完全交叉である。
実際単元三次小行列で核を解けばその方程式と残った一つの行列式とに局所分解する。
グラフの積空間内の類は $(l+h)^4$ であり、これから $lh=6$、他側も同様である。
余因子三次式は同時に消えず、左右核のSegre座標を与えるので $l+r=3h$。
完全交叉のKoszul列を $(1,1)$ 捩れば、十六双線型切断から $H^0(l+h)$ への制限は全射で、
その核は $t$ の四縮約である。反対も同様である。
実際負の中間項には $\mathcal O(-j,-j)$、$1\le j\le3$ しか出ず非零コホモロジーを持たない。
六関係中の残り二式が鉛筆を与え、元の六式の全零点はまさにその基底スキームである。
中央射影が点グラフの両側を同じスキームに同定するので二鉛筆の基底は一致する。
$A^2=B^2=20, AB=30$ である。

逆に両方の基底が有限ならそれぞれ長さ20であり、後者が前者を含むので等しい。
有限基底 $X$ の完全交叉列 $0\to\mathcal O(-B)\to\mathcal O^2\to I_{X/T}(B)\to0$
から $H^0(I_{X/T}(B))$ は鉛筆そのものである（$B$ は豊富で $H^1(-B)=0$）。
元の二窓の六式は各鉛筆の四核つき逆像に等しい。
$t$ は三次重なりに入り、系\ref{cor:two-window-extension}の復元によりこれがちょうど印付き四窓を与える。
最後の有限検査は正則性を仮定せずに復元テンソルに施せる。基底変更と逆操作が両立するので同型条件も従う。
\end{proof}

\begin{proposition}\label{prop:finite-marked-eight-determinant}
$\mathrm{NS}(T)=\mathbb Z l\oplus\mathbb Z h$ なら、有限基底を持つ上記の鉛筆で
反対鉛筆の候補は八次元空間間の乗法だけから決まる。
\[
 H^0(I_{X/T}(B))=\ker\big[H^1(B-2A)\longrightarrow
                 \Lambda^*\otimes H^1(B-A)\big]
\]
と体積因子を除き同定でき、両 $H^1$ は八次元である。
従って必要な第一の退化条件は八列・十六行の行列が階数6であることとなる。
\end{proposition}
\begin{proof}
$B-A=3h-2l, B-2A=2h-3l$ はともに平方 $-20$。
格子の平方は4の倍数なので曲面には $(-2)$ 曲線がない。
従ってこれらの因子とその負は有効でなく、Riemann--Rochで $h^1=8$ を得る。
二切断のKoszul列を $B$ 捩ると主張の核になる。
行列の各ブロックは対応する $A$ の切断による乗法である。
\end{proof}

\begin{theorem}\label{thm:dominant-four-by-four-by-four-slice}
十五次元混合成分に一般の $x$ を印し三文字へ縮約すると、
\emph{一般の全ての}四行四列四層テンソルの射影同値類に達し、
その印付き拡張はこの成分で生成的有限である。
従って前定理の滑らかなK3切断は十五次元成分の稠密開集合を被覆する。
\end{theorem}
\begin{proof}
定理\ref{thm:full-fifteen-component}の有理点 $t_0$ で $x=(1,2,3,5)$ を固定する。
基底つき二切断階数6軌跡はそこで滑らかで次元76である。
各切断行列の左右核から $u^t\dot w v=0$ の百式を作り、二組を同じ256座標に戻す。
素数1009でこの連立核は次元76、$\dot w\mapsto\dot w(x)$ の核への制限は階数64である。
元の二切断の核は局所整数環で直和因子であり、連立核次元が標数零と一致するので接基底が持ち上がる。
従って標数零でこの縮約は支配的である。
同じ点で三方向の行列式四次式の四偏微分を全六次単項式倍すると、
三方向とも九次全空間の次元220を張る。従って三つの射影曲面は滑らかである。
正確な行列・座標順と剰余検算は \path{computations/verify_determinantal_marked_slice.py} にある。
一般のファイバーの基底付き次元は $76-64=12$ であり、
$x$ を線型的に固定する第一空間の基底変更が既に次元12を占める。
実際四テンソルの一文字支持は充満なので、それだけを動かしてテンソルを固定する元は恒等である。
固定した $x$ を任意の印へ移すこと、および三文字の基底変更を両側で割ることから結論が従う。
\end{proof}
